%----------------------------------------------------------------------------------------
% BETROUWBAARHEIDSINTERVALLEN. Alle opgaven zijn eigen werk.
%
% Opbouw: Start, leerstof (theorie, werkwijze of uitleg, voorbeeld, korte oefening), en op een
% nieuwe pagina de Opdrachten (door elkaar, stijgende moeilijkheid), Koppel en Test jezelf.
% Labels van dit hoofdstuk beginnen met h4-.
%----------------------------------------------------------------------------------------

\mychapter{4}{Betrouwbaarheidsintervallen}

\lkr{Leerplan B+S'': keuzedoel K1 (de betekenis van betrouwbaarheidsniveau, betrouwbaarheidsinterval en foutenmarge uitleggen in betekenisvolle situaties), met de wenken: het effect van niveau, steekproefgrootte en standaardafwijking op de foutenmarge, intervallen voor proporties en gemiddelden, en ICT. Dit hoofdstuk is facultatief. Op de TI-84: \texttt{1-PropZInt} en \texttt{ZInterval} in het menu STAT, TESTS.}

Een krant meldt: \textit{23~\% van de Belgen overweegt een elektrische auto, met een foutenmarge van 2,6~procentpunt}. Wat betekent die foutenmarge? Een steekproef geeft nooit precies de parameter, maar je kunt wel zeggen in welk interval die waarschijnlijk ligt. In dit hoofdstuk leer je zo'n interval berekenen, en vooral: het juist interpreteren.

\begin{oefeningen}[Start]
\oef{h4-start-1}

\oef{h4-verken}
\end{oefeningen}

%========================================================================================
\section{Een interval voor een proportie}

\subsection{Het 95~\%-betrouwbaarheidsinterval}

Bij een aselecte steekproef van grootte $n$ is $\hat{p}$ bij benadering normaal verdeeld met gemiddelde $p$ en standaardafwijking $\sigma(\hat{p}) = \sqrt{\frac{p(1 - p)}{n}}$. Omdat $P(-1{,}96 \leqslant Z \leqslant 1{,}96) \approx 0{,}95$, ligt $\hat{p}$ in 95~\% van de steekproeven binnen $1{,}96\,\sigma(\hat{p})$ van $p$. Omgekeerd ligt $p$ dan binnen die afstand van $\hat{p}$. Omdat je $p$ niet kent, vervang je het in de formule voor $\sigma(\hat{p})$ door $\hat{p}$.

\begin{theorie}
Een \textbf{\textcolor{accent}{95~\%-betrouwbaarheidsinterval}} voor een populatieproportie $p$, op basis van een aselecte steekproef met aandeel $\hat{p}$, is

\vspace{1mm}
\qquad $\formule{\Bigl[\hat{p} - E,\ \hat{p} + E\Bigr]}$ \quad met \textbf{\textcolor{accent}{foutenmarge}} \quad $\formule{E = 1{,}96 \sqrt{\dfrac{\hat{p}(1 - \hat{p})}{n}}}$

\vspace{1mm}
Het getal 95~\% is het \textbf{\textcolor{accent}{betrouwbaarheidsniveau}}.
\end{theorie}

\voorbeeld{}{In een aselecte steekproef van 1000 Belgen overweegt 230 een elektrische auto. Bereken een 95~\%-betrouwbaarheidsinterval voor het aandeel van alle Belgen.\label{vb:h4-proportie}}{%
\uitleg{We berekenen $\hat{p}$ en de foutenmarge, en controleren met het rekentoestel.}
\begin{stappenbreed}
\hat{p} = \frac{230}{1000} = 0{,}23 & Het aandeel in de steekproef. \\[3ex]
E = 1{,}96\sqrt{\frac{0{,}23 \cdot 0{,}77}{1000}} \approx 0{,}026 & De foutenmarge: 2,6 procentpunt. \\[3ex]
[0{,}204;\ 0{,}256] & Met \texttt{1-PropZInt}: hetzelfde interval.
\end{stappenbreed}
Met 95~\% betrouwbaarheid overweegt tussen 20,4~\% en 25,6~\% van alle Belgen een elektrische auto.}

\begin{werkwijze}
Bereken de schatting $\hat{p}$ (of $\xbar$) en de foutenmarge $E$. Het interval is de schatting plus en min $E$. Formuleer het besluit over de parameter van de hele populatie, met het betrouwbaarheidsniveau.
\end{werkwijze}

\begin{oefeningen}
\oef{h4-proportie}
\end{oefeningen}

De formule steunt op twee voorwaarden. De steekproef moet aselect zijn: een interval rond een vertekende schatting zegt niets over $p$. En de normale benadering moet bruikbaar zijn, zoals in hoofdstuk 3: $n\hat{p} \geqslant 5$ en $n(1 - \hat{p}) \geqslant 5$. Hebben bijvoorbeeld 2 van 40 personen een eigenschap, dan is $\hat{p} = 0{,}05$ en $E = 1{,}96\sqrt{\frac{0{,}05 \cdot 0{,}95}{40}} \approx 0{,}068$. Het interval $[-0{,}018;\ 0{,}118]$ bevat dan negatieve aandelen, wat onmogelijk is: met $n\hat{p} = 2$ mag je de formule niet gebruiken.

\begin{oefeningen}
\oef{h4-voorwaarde}
\end{oefeningen}

\subsection{Wat betekent 95~\%?}

Het interval hangt af van de steekproef: een andere steekproef geeft een ander interval. Wat is dan de betekenis van die 95~\%? Ze gaat over de methode, niet over één interval.

\begin{theorie}
\begin{stelling}{Betrouwbaarheid van de methode}
Voor grote $n$ bevat het interval $[\hat{p} - E,\ \hat{p} + E]$ de echte $p$ bij ongeveer 95~\% van alle aselecte steekproeven:

\vspace{1mm}
\qquad $\formule{P(\hat{p} - E \leqslant p \leqslant \hat{p} + E) \approx 0{,}95}$
\end{stelling}

\begin{bewijs}
Tel bij de linkerongelijkheid $E$ op en trek van de rechterongelijkheid $E$ af:
\bewijsregel{\hat{p} - E \leqslant p \leqslant \hat{p} + E \iff p - E \leqslant \hat{p} \leqslant p + E .}
Het interval bevat $p$ dus precies als $\hat{p}$ op afstand hoogstens $E$ van $p$ ligt. Voor grote $n$ ligt $\hat{p}$ dicht bij $p$, zodat $E \approx 1{,}96\,\sigma(\hat{p})$, en is $\hat{p}$ bij benadering $N(p, \sigma(\hat{p}))$. Standaardiseren geeft
\bewijsregel{P\bigl(p - 1{,}96\,\sigma(\hat{p}) \leqslant \hat{p} \leqslant p + 1{,}96\,\sigma(\hat{p})\bigr) = P(-1{,}96 \leqslant Z \leqslant 1{,}96) \approx 0{,}95 .}
\end{bewijs}
\end{theorie}

\naast{De twee benaderingen in het bewijs, de normale verdeling en $\hat{p}$ in plaats van $p$, maken van de 95~\% een \textit{ongeveer}. De figuur toont een simulatie: 20 aselecte steekproeven van 200 personen uit een populatie met $p = 0{,}4$, met telkens het 95~\%-interval. Negentien intervallen bevatten 0,4, één niet. Voor één bepaald interval weet je niet of het bij de gelukkige 95~\% hoort.}{\begin{tikzpicture}[x=16cm,y=.24cm,font=\scriptsize]
  \draw[grijs] (.27,0) -- (.59,0);
  \foreach \t/\l in {.3/{0,3},.4/{0,4},.5/{0,5}}{\draw[grijs] (\t,0) -- (\t,-.4) node[below,text=black] {\l};}
  \draw[fel!80!black,dashed] (.4,0) -- (.4,21) node[above,text=black] {$p$};
  \foreach \a/\b/\i in {.289/.421/1,.352/.488/2,.322/.458/4,.366/.504/5,.327/.463/6,.342/.478/7,.371/.509/8,.303/.437/9,.298/.432/10,.337/.473/11,.318/.452/12,.293/.427/13,.366/.504/14,.337/.473/15,.303/.437/16,.361/.499/17,.293/.427/18,.313/.447/19,.298/.432/20}{\draw[accent,thick] (\a,\i) -- (\b,\i); \fill[accent] ({(\a+\b)/2},\i) circle (1pt);}
  \draw[fel!80!black,very thick] (.426,3) -- (.564,3); \fill[fel!80!black] (.495,3) circle (1.2pt);
\end{tikzpicture}}

\voorbeeld{}{Na het interval $[0{,}204;\ 0{,}256]$ uit voorbeeld \ref{vb:h4-proportie} schrijft een krant: \textit{De kans is 95~\% dat het aandeel tussen 20,4~\% en 25,6~\% ligt.} Een lezer schrijft: \textit{95~\% van de Belgen ligt in dat interval.} Wat is er fout? Hoe zeg je het juist?\label{vb:h4-interpretatie}}{%
\uitleg{We gaan na wat vastligt en wat van het toeval afhangt, en waarover het interval gaat.}
De krant vergist zich: $p$ ligt vast, en ligt in dit interval of niet. Het toeval zat in de steekproef, en die is al getrokken. De lezer vergist zich ook: het interval gaat over één parameter, het aandeel $p$, niet over personen. Juist is: \textit{met een methode die in 95~\% van de steekproeven een interval rond de echte $p$ geeft, vonden we $[0{,}204;\ 0{,}256]$.}}

\begin{oefeningen}
\oef{h4-interpretatie}
\end{oefeningen}

\subsection{Een ander betrouwbaarheidsniveau}

Het getal 1,96 komt van 95~\%: tussen $-1{,}96$ en $1{,}96$ ligt 95~\% van de standaardnormale verdeling, dus links van 1,96 ligt 97,5~\%. Wil je vaker gelijk hebben, dan heb je een groter getal $z$ nodig, en een breder interval.

\begin{theorie}
Bij een betrouwbaarheidsniveau van 90~\%, 95~\% of 99~\% is de foutenmarge \quad $\formule{E = z \sqrt{\dfrac{\hat{p}(1 - \hat{p})}{n}}}$ \quad met

\vspace{1mm}
{\centering\arrayrulecolor{accent}\renewcommand{\arraystretch}{1.4}%
\begin{tabular}{|c|c|c|c|}
\hline
niveau & 90~\% & 95~\% & 99~\% \\ \hline
$z$ & $\formule{1{,}645}$ & $\formule{1{,}96}$ & $\formule{2{,}576}$ \\ \hline
\end{tabular}\par}
\end{theorie}

Je vindt elke $z$ zoals 1,96: voor 90~\% ligt links van $z$ 95~\%, dus $z = {}$\texttt{invNorm(0.95,0,1)} $\approx 1{,}645$.

\voorbeeld{}{Bereken voor de 230 van de 1000 Belgen uit voorbeeld \ref{vb:h4-proportie} een 99~\%-\hspace{0pt}betrouwbaarheidsinterval. Vergelijk met het 95~\%-interval.\label{vb:h4-niveau}}{%
\uitleg{Alleen $z$ verandert: 2,576 in plaats van 1,96.}
\begin{stappenbreed}
E = 2{,}576\sqrt{\frac{0{,}23 \cdot 0{,}77}{1000}} \approx 0{,}034 & Groter dan 0,026. \\[3ex]
[0{,}196;\ 0{,}264] & Met \texttt{1-PropZInt}, met niveau 0.99.
\end{stappenbreed}
Het 99~\%-interval is breder: meer zekerheid kost nauwkeurigheid.}

\begin{oefeningen}
\oef{h4-niveau}
\end{oefeningen}

%========================================================================================
\section{Een interval voor een gemiddelde}

Voor een gemiddelde gebruik je de steekproevenverdeling van $\overline{X}$, met standaardafwijking $\frac{\sigma}{\sqrt n}$. Ook hier kennen we in deze cursus $\sigma$.

\begin{theorie}
Een \textbf{\textcolor{accent}{95~\%-betrouwbaarheidsinterval}} voor een populatiegemiddelde $\mu$, met $\sigma$ gekend, op basis van een aselecte steekproef met gemiddelde $\xbar$, is

\vspace{1mm}
\qquad $\formule{\Bigl[\xbar - E,\ \xbar + E\Bigr]}$ \quad met \quad $\formule{E = 1{,}96 \dfrac{\sigma}{\sqrt{n}}}$
\end{theorie}

Dat dit interval $\mu$ bevat bij 95~\% van de steekproeven, toon je zoals bij een proportie, en eenvoudiger: $\sigma$ is gekend (oefening \ref{oef:h4-o20}). Voor 90~\% of 99~\% vervang je 1,96 door $z$ uit de tabel.

\voorbeeld{}{Een fabriek test 50 batterijen. Ze gaan gemiddeld 31 uur mee. Uit ervaring is $\sigma = 4$ uur. Bereken een 95~\%-betrouwbaarheidsinterval voor de gemiddelde levensduur van alle batterijen.\label{vb:h4-gemiddelde}}{%
\uitleg{De foutenmarge is 1,96 standaardafwijkingen van het steekproefgemiddelde.}
\begin{stappenbreed}
E = 1{,}96 \cdot \frac{4}{\sqrt{50}} \approx 1{,}11 & De foutenmarge in uur. \\[1ex]
[29{,}9;\ 32{,}1] & Met \texttt{ZInterval}: hetzelfde.
\end{stappenbreed}
De gemiddelde levensduur van alle batterijen ligt met 95~\% betrouwbaarheid tussen 29,9 en 32,1 uur.}

\begin{oefeningen}
\oef{h4-gemiddelde}
\end{oefeningen}

Een interval ligt symmetrisch rond de schatting. Krijg je alleen het interval, dan kun je dus terugrekenen.

\voorbeeld{}{Een onderzoeker meldt het 95~\%-betrouwbaarheidsinterval $[82{,}06;\ 87{,}94]$ voor een gemiddelde, met $\sigma = 12$. Wat was het steekproefgemiddelde, de foutenmarge en de steekproefgrootte?\label{vb:h4-terug}}{%
\uitleg{Het midden van het interval is $\xbar$, de halve breedte is $E$. Uit $E$ halen we daarna $n$.}
\begin{stappen}
\xbar = (82{,}06 + 87{,}94) : 2 = 85 & Het midden. \\[.6ex]
E = (87{,}94 - 82{,}06) : 2 = 2{,}94 & De halve breedte. \\[.6ex]
1{,}96 \cdot \frac{12}{\sqrt n} = 2{,}94 \ \Rightarrow\ \sqrt n = 8 & Dus $n = 64$.
\end{stappen}}

\begin{oefeningen}
\oef{h4-terug}
\end{oefeningen}

%========================================================================================
\section{De foutenmarge beïnvloeden}

\subsection{De steekproefgrootte}

De foutenmarge $E = z \cdot \frac{\sigma}{\sqrt n}$ (of met $\sqrt{\hat{p}(1 - \hat{p})}$ in de plaats van $\sigma$) hangt af van drie dingen: het betrouwbaarheidsniveau via $z$, dat zag je al, de spreiding en de steekproefgrootte. Een grotere steekproef geeft een kleinere marge, maar trager dan je zou denken.

\begin{theorie}
\begin{stelling}{Vier keer zoveel waarnemingen}
Maak je de steekproef $k^2$ keer zo groot, dan wordt de foutenmarge $k$ keer zo klein:

\vspace{1mm}
\qquad $\formule{z \dfrac{\sigma}{\sqrt{k^2 n}} = \dfrac{1}{k} \cdot z \dfrac{\sigma}{\sqrt{n}}}$
\end{stelling}

\begin{bewijs}
$\sqrt{k^2 n} = k\sqrt{n}$ voor $k > 0$. Hetzelfde geldt met $\sqrt{\hat{p}(1 - \hat{p})}$ in de plaats van $\sigma$, als $\hat{p}$ ongeveer gelijk blijft.
\end{bewijs}
\end{theorie}

Vier keer zoveel waarnemingen halveert dus de marge, en voor een tien keer kleinere marge heb je honderd keer zoveel waarnemingen nodig. Nauwkeurigheid is duur.

\voorbeeld{}{De peiling bij 1000 Belgen uit voorbeeld \ref{vb:h4-proportie} had een foutenmarge van 2,6 procentpunt. Hoeveel personen moet het bureau ondervragen om die marge te halveren, als $\hat{p}$ ongeveer 0,23 blijft?\label{vb:h4-wortel}}{%
\uitleg{De marge halveren is $k = 2$: dan moet $n$ maal $k^2 = 4$.}
\begin{stappen}
n = 4 \cdot 1000 = 4000 & Vier keer zoveel personen. \\[1ex]
E = 1{,}96\sqrt{\tfrac{0{,}23 \cdot 0{,}77}{4000}} \approx 0{,}013 & Controle: 1,3 procentpunt.
\end{stappen}}

\begin{oefeningen}
\oef{h4-marge}
\end{oefeningen}

\subsection{De spreiding}

Hoe groter de spreiding $\sigma$, hoe breder het interval, en wel evenredig: een populatie met dubbel zo grote $\sigma$ geeft bij dezelfde $n$ een dubbel zo grote marge. Bij een proportie speelt $\sqrt{\hat{p}(1 - \hat{p})}$ de rol van $\sigma$, en die hangt van $\hat{p}$ zelf af.

\begin{theorie}
\begin{stelling}{Grootste spreiding bij een half}
Voor elke $p$ is \quad $\formule{p(1 - p) \leqslant \dfrac14}$, \quad met gelijkheid alleen voor $p = \frac12$.
\end{stelling}

\begin{bewijs}
Werk uit en splits een kwadraat af:
\bewijsregel{\dfrac14 - p(1 - p) = p^2 - p + \dfrac14 = \Bigl(p - \dfrac12\Bigr)^2 \geqslant 0 ,}
met gelijkheid alleen als $p - \frac12 = 0$.
\end{bewijs}
\end{theorie}

Een aandeel rond 50~\% is dus het moeilijkst te schatten. Wil je vooraf de steekproefgrootte bepalen, en ken je $\hat{p}$ nog niet, dan reken je met het ongunstigste geval $\hat{p} = 0{,}5$: dan is je marge zeker niet groter dan gepland.

\subsection{De nodige steekproefgrootte}

Met die twee stellingen kun je vooraf berekenen hoeveel waarnemingen je nodig hebt voor een gewenste foutenmarge.

\voorbeeld{}{Een peilingbureau wil een aandeel schatten met een foutenmarge van hoogstens 3 procentpunt, bij 95~\% betrouwbaarheid. Hoeveel personen moet het minstens ondervragen?\label{vb:h4-steekproefgrootte}}{%
\uitleg{We lossen $1{,}96\sqrt{\frac{0{,}25}{n}} \leqslant 0{,}03$ op naar $n$, in het ongunstigste geval.}
\begin{stappen}
\sqrt{n} \geqslant \frac{1{,}96 \cdot 0{,}5}{0{,}03} \approx 32{,}7 & Want $\sqrt{0{,}25} = 0{,}5$. \\[3ex]
n \geqslant 1067{,}1 \ \Rightarrow\ n = 1068 & Naar boven afronden. Daarom ondervragen peilingen vaak ongeveer 1000 personen.
\end{stappen}}

\begin{werkwijze}
Schrijf de foutenmarge als functie van $n$ en eis dat ze hoogstens de gewenste marge is. Neem bij een proportie $\hat{p} = 0{,}5$, tenzij je een betere schatting hebt. Los op naar $n$ en rond altijd naar boven af.
\end{werkwijze}

\begin{oefeningen}
\oef{h4-n}
\end{oefeningen}

%========================================================================================
\opdrachtenkop

\begin{oefeningen}
\needspace{9\baselineskip}
\oef{h4-d1}

\oef{h4-o2}

\needspace{14\baselineskip}
\oef{h4-o11}

\needspace{10\baselineskip}
\oef{h4-o9}

\needspace{10\baselineskip}
\oef{h4-d2}

\needspace{10\baselineskip}
\oef{h4-o10}

\oef{h4-o1}

\needspace{14\baselineskip}
\oef{h4-o8}

\oef{h4-o12}

\needspace{10\baselineskip}
\oef{h4-o3}

\needspace{16\baselineskip}
\oef{h4-d3}

\needspace{12\baselineskip}
\oef{h4-o15}

\needspace{16\baselineskip}
\oef{h4-o13}

\needspace{12\baselineskip}
\oef{h4-o19}

\needspace{12\baselineskip}
\oef{h4-o14}

\oef{h4-o4}

\needspace{16\baselineskip}
\oef{h4-o16}

\oef{h4-o5}

\needspace{12\baselineskip}
\oef{h4-o17}

\needspace{12\baselineskip}
\oef{h4-o18}

\oef{h4-o6}

\needspace{14\baselineskip}
\oef{h4-o20}

\needspace{12\baselineskip}
\oef{h4-o23}

\needspace{12\baselineskip}
\oef{h4-o21}

\oef{h4-o7}

\needspace{12\baselineskip}
\oef{h4-o22}
\end{oefeningen}

%----------------------------------------------------------------------------------------
% KOPPEL: verbanden met andere leerstof
%----------------------------------------------------------------------------------------
\needspace{20\baselineskip}
\opdrachtsectie{Koppel}

\lkr{De eerste oefening gebruikt het verloop van een tweedegraadsfunctie (2de graad), de tweede de binomiale verdeling.}

\begin{oefeningen}
\oef{h4-koppel-1}

\oef{h4-koppel-2}
\end{oefeningen}

%----------------------------------------------------------------------------------------
\opdrachtsectie{Test jezelf}

\begin{oefeningen}
\oef{h4-t1}

\oef{h4-t2}

\oef{h4-t3}

\oef{h4-t4}

\oef{h4-t5}

\oef{h4-t6}

\oef{h4-t7}
\end{oefeningen}

\herhaalschema{%
  \fv{oef:h4-t1}{vb:h4-proportie}{oef:h4-proportie}%
  \fv{oef:h4-t2}{vb:h4-gemiddelde}{oef:h4-gemiddelde}%
  \fv{oef:h4-t3}{vb:h4-steekproefgrootte}{oef:h4-n}%
  \fv{oef:h4-t4}{vb:h4-interpretatie}{oef:h4-interpretatie}%
  \fv{oef:h4-t5}{vb:h4-niveau}{oef:h4-niveau}%
  \fv{oef:h4-t6}{vb:h4-terug}{oef:h4-terug}%
  \fv{oef:h4-t7}{vb:h4-wortel}{oef:h4-marge}}

\ifzelfstudie
\section*{Maak je eigen samenvatting}
Sluit dit boek en schrijf in je schrift: hoe je een 95~\%-betrouwbaarheidsinterval voor een proportie en voor een gemiddelde berekent, wat de foutenmarge en het betrouwbaarheidsniveau betekenen, de juiste interpretatie van \textit{95~\% betrouwbaar} (en twee foute), en hoe niveau, steekproefgrootte en spreiding de foutenmarge beïnvloeden. Vergelijk daarna met de kaders.
\fi
